% Created 2022-10-06 Thu 17:59
% Intended LaTeX compiler: pdflatex
\documentclass[12pt,oneside,a4paper,english]{abntex2}
\usepackage[utf8]{inputenc}
\usepackage[T1]{fontenc}
\usepackage{graphicx}
\usepackage{longtable}
\usepackage{wrapfig}
\usepackage{rotating}
\usepackage[normalem]{ulem}
\usepackage{amsmath}
\usepackage{amssymb}
\usepackage{capt-of}
\usepackage{hyperref}
\usepackage{lmodern}
\usepackage[T1]{fontenc}
\usepackage[utf8]{inputenc}
\usepackage{indentfirst}
\usepackage{color}
\usepackage{graphicx}
\usepackage{microtype}
%\usepackage{cmap}
\usepackage{lastpage}
\usepackage{epstopdf}
\usepackage{lipsum}
\usepackage{nomencl}
\usepackage{amsmath}
\usepackage{bbm}
\usepackage{multirow}
\usepackage{rotating}
\usepackage{pdfpages}
\usepackage[font=footnotesize]{subfig}
\usepackage{booktabs}
\usepackage{pdflscape}
\usepackage{chngcntr}
\let\printglossary\relax
\let\theglossary\relax
\let\endtheglossary\relax
\usepackage[nonumberlist,acronym,nomain]{glossaries}
\newglossary[tlg]{simbolos}{tld}{tdn}{Lista de símbolos}
\makeglossaries
\counterwithin{figure}{chapter}
\counterwithin{table}{chapter}
\usepackage[brazilian,hyperpageref]{backref}
\usepackage[alf]{abntex2cite}
\usepackage{unicamp}
\renewcommand{\backrefpagesname}{Citado na(s) p\'{a}gina(s):~}
\renewcommand{\backref}{}
\definecolor{blue}{RGB}{41,5,195}
\makeatletter
\usepackage{hyperref}
\makeatother
\makeatletter
\setlength{\@fptop}{5pt} % Set distance from top of page to first float
\makeatother
\newcommand{\quadroname}{Quadro}
\newcommand{\listofquadrosname}{Lista de quadros}
\newfloat[chapter]{quadro}{loq}{\quadroname}
\newlistof{listofquadros}{loq}{\listofquadrosname}
\newlistentry{quadro}{loq}{0}
\setfloatadjustment{quadro}{\centering}
\counterwithout{quadro}{chapter}
\renewcommand{\cftquadroname}{\quadroname\space}
\renewcommand*{\cftquadroaftersnum}{\hfill--\hfill}
\setfloatlocations{quadro}{hbtp}
\setlength{\parindent}{2cm}
\titulo{Dissertation/Thesis Title}
\autor{Author Name}
\local{Campinas}
\data{Year}
\orientador{Prof. Dr. Supervisor}
\coorientador[Coorientador]{Prof. Dr. Co-supervisor}
\instituicao{UNIVERSIDADE ESTADUAL DE CAMPINAS \par Faculdade de Engenharia Elétrica e de Computação}
\tipotrabalho{Thesis (Doctorate)}
\preambulo{Thesis presented to the School of Electrical and Computer Engineering of the University of Campinas in partial fulfillment of the requirements for the degree of Doctor in Electrical Engineering, in the area of Computer Engineering.}
\makeindex
\renewcommand\maketitle{}
\date{Year}
\title{Dissertation/Thesis Title}
\hypersetup{
 pdfauthor={Author Name},
 pdftitle={Dissertation/Thesis Title},
 pdfkeywords={},
 pdfsubject={},
 pdfcreator={Emacs 28.2 (Org mode 9.6)}, 
 pdflang={English}}
\begin{document}

\maketitle
\imprimircapa


\imprimirfolhaderosto*


\begin{fichacatalografica}
    \vspace*{\fill}
    \begin{center}
        \textsc{Inclua aqui o pdf com a ficha catalogr\'{a}fica fornecida pela BAE.}
    \end{center}
    \vspace*{\fill}
    %\includepdf{fig_ficha_catalografica.pdf}
\end{fichacatalografica}


\begin{folhadeaprovacao}
  \begin{center}
    COMISS\~{A}O JULGADORA - TESE DE DOUTORADO
    %\textsc{Inclua aqui a folha de assinaturas.}
\end{center}
\noindent
\begin{minipage}{\textwidth}\SingleSpacing
Candidato(a): Nome do Autor      RA: XXXXXX
Data de defesa: XX de MES de 202X
Título da Tese: "XXXXXXXXXXXXXXXXXXXXXXXXXXXXXXX"
\vspace{2cm}

Profa. Dra. Xxxxxxxxxx (Presidente)

Profa. Dra. xxxxxxx

Profa. Dra. xxxxxxx

Profa. Dra. xxxxxxxxx

Profa. Dra xxxxxxxxxxxx

\vspace{2cm}

A Ata de Defesa, com as respectivas assinaturas dos membros da Comiss\~{a}o Julgadora, encontra-se no SIGA (Sistema de Fluxo de Disserta\c{c}\~{a}o/Tese) e na Secretaria de P\'{o}s-Gradua\c{c}\~{a}o da Faculdade de Engenharia El\'{e}trica e de Computa\c{c}\~{a}o.
\end{minipage}

\end{folhadeaprovacao}


\begin{dedicatoria}
   \vspace*{\fill}
   \centering
   \noindent
   \textit{ Dedico esta tese \`{a} todo mundo.} \vspace*{\fill}
\end{dedicatoria}


\begin{agradecimentos}
    Escreva seus agradecimentos.
    Obs.: Ser\'{a} obrigat\'{o}rio caso o autor tenha recebido aux\'{\i}lio financeiro, parcial ou integral, de ag\^{e}ncia (s) de fomento. Neste caso, ele dever\'{a} atender \`{a} legisla\c{c}\~{a}o vigente espec\'{\i}fica de cada uma das Ag\^{e}ncias quanto \`{a} necessidade de se fazer refer\^{e}ncia ao apoio recebido e ao n\'{u}mero de processo. (OF PRPG 002/2019 – Orienta\c{c}\~{a}o sobre disserta\c{c}\~{o}es e teses).

    -	Em caso de Bolsa CAPES, usar a seguinte express\~{a}o, como indicado na portaria 206 da CAPES:

    O presente trabalho foi realizado com apoio da Coordena\c{c}\~{a}o de Aperfei\c{c}oamento de Pessoal de N\'{\i}vel Superior - Brasil (CAPES) - C\'{o}digo de Financiamento 001.

    -	Em caso de bolsa CNPQ, usar a seguinte express\~{a}o de agradecimento:

    O presente trabalho foi realizado com apoio do CNPq, Conselho Nacional de Desenvolvimento Cient\'{\i}fico  e Tecnol\'{o}gico – Brasil.

    -	Em caso de bolsa FAPESP, fazer agradecimento contendo nome FAPESP, o n\'{u}mero do processo FAPESP a que se refere este Termo de Outorga, no modelo:

    processo nº aaaa/nnnnn-d, Funda\c{c}\~{a}o de Amparo \`{a} Pesquisa do Estado de S\~{a}o Paulo (FAPESP).

    Os artigos escritos em idioma estrangeiro dever\~{a}o indicar o apoio da FAPESP em ingl\^{e}s, conforme o seguinte modelo: grant \# aaaa/nnnnn-d, S\~{a}o Paulo Research Foundation (FAPESP).
\end{agradecimentos}


\begin{epigrafe}
    \vspace*{\fill}
    \begin{flushright}
        \textit{``Escreva aqui a sua ep\'{\i}grafe (Opcional)''\\
        (Cita\c{c}\~{a}o)}
    \end{flushright}
\end{epigrafe}


\setlength{\absparsep}{18pt}
\begin{resumo}
 Insira seu resumo. (Obrigat\'{o}rio, em portugu\^{e}s m\'{a}ximo de 500 palavras)

    \lipsum[1]

    \vspace{\onelineskip}

    \noindent\textbf{Palavras-chaves}: palavra-chave 1; palavra-chave 2; palavra-chave 3.
\end{resumo}


\begin{resumo}[Abstract]
 \begin{otherlanguage*}{english}
   Same content of "Resumo".

   \vspace{\onelineskip}

   \noindent
   \textbf{Keywords}: latex. abntex. text editoration.
 \end{otherlanguage*}
\end{resumo}


\pdfbookmark[0]{\listfigurename}{lof}
\listoffigures*
\cleardoublepage


\pdfbookmark[0]{\listtablename}{lot}
\listoftables*
\cleardoublepage


\newacronym{NIST}{NIST}{National Institute of Standards and Technology}
\newacronym{TWC}{TWC}{Transformada Wavelet Contínua}
\newacronym{TWD}{TWD}{Transformada Wavelet Discreta}
\printglossary[type=\acronymtype,title=\{List of Abbreviations\}]
\cleardoublepage

\printglossary[type=simbolos,title=\{List of Symbols\}]
\cleardoublepage


\pdfbookmark[0]{\contentsname}{toc}
\tableofcontents*
\cleardoublepage


\textual

bold:
\textbf{surround text with star}

italic:
\emph{surround text with slash}

underline:\#
\uline{surround \# text with underscore}

block code\# : (<s)
\#+BEGIN\textsubscript{SR}\# C python
print(``Hello'')
\#+END\textsubscript{SRC}

links: \begin{center}
\includegraphics[width=.9\linewidth]{./figures/unicamp4.png}
\end{center}

M-x org-toggle-inline-images (C-c C-x C-v)

\section{Introduction}
\label{sec:org43e0daa}
\subsection{Part I}
\label{sec:orgb22c524}
test 1
test 2
test 3

\subsection{Part II}
\label{sec:orgc299b81}
test 4
test 5
test 6

\section{Methods}
\label{sec:org8daffab}
asdfasdfasdf
\end{document}
